%% kapa: exam=1 round=30 year=2019 period=1 form=A subject=civil_law subject_ko=민법 seq=11 number=11
%% source: Q-Net 감정평가사 기출문제 (cst00309) artlSeq=5205205 file=extracted/1008/5205205/2019년 제30회 감정평가사 1차시험 문제지/2019년 제30회 감정평가사 1차시험 문제지/2019년도 제30회 감정평가사 1차 1교시 A형.hwp
%% conversion: kordoc 4.24.51 -> tools/gen1.py (auto); review=proofread-sonnet; flags=none
\begin{question}{11}
통정허위표시에 관한 설명으로 옳지 않은 것은? (다툼이 있으면 판례에 따름)
\begin{choices}
  \item 상대방과 통정한 허위의 의사표시는 무효이지만, 이러한 무효는 과실로 인하여 허위표시라는 사실을 인식하지 못한 제3자에게 대항할 수 없다.
  \item 강제집행을 면할 목적으로 부동산에 허위의 근저당권설정등기를 경료하는 행위는 민법 제103조의 선량한 풍속 기타 사회질서에 위반한 사항을 내용으로 하는 법률행위이다.
  \item 선의의 제3자에 대해서는 통정허위표시의 당사자뿐만 아니라 그 누구도 통정허위표시의 무효로 대항할 수 없다.
  \item 부동산의 가장양수인으로부터 해당 부동산을 취득한 제3자 A가 악의이고, 그로부터 그 부동산을 전득한 B가 선의라면 통정허위표시의 무효로써 B에게 대항할 수 없다.
  \item 당사자들이 실제로는 증여계약을 체결하면서 매매계약인 것처럼 통정허위표시를 하였다면 은닉행위인 증여계약은 유효할 수 있다.
\end{choices}
\end{question}
\answer{2}
